\chapter{Testowanie i ocena opracowanego rozwiązania}
\label{chap:testowanie}

Ocena systemu PlanWise wymaga rozróżnienia poprawności działania aplikacji od użyteczności generowanych rekomendacji. Poprawne zapisanie zadania, wyznaczenie wartości funkcji oceny lub odebranie odpowiedzi modelu językowego nie oznacza jeszcze, że system skutecznie wspomaga zarządzanie projektem. Konieczne jest zatem sprawdzenie zarówno realizacji przypadków użycia, jak i jakości mechanizmów analitycznych.

W rozdziale przedstawiono zakres testów obecnych w repozytorium, scenariusze weryfikacji funkcjonalnej oraz metodykę eksperymentów dotyczących ryzyka, priorytetów, przydziałów i kosztów. Oddzielono ustalenia wynikające z analizy kodu od wyników wymagających uruchomienia aplikacji i przeprowadzenia pomiarów.

Zakres przedstawionych wyników jest przy tym nierówny i zostało to w rozdziale jawnie zaznaczone. Dysponujemy raportem z wykonania zautomatyzowanego zestawu testów oraz wynikami eksperymentów dotyczących modelu predykcji opóźnień, przeprowadzonych na zewnętrznych zbiorach danych. Nie przeprowadzono natomiast eksperymentów porównawczych dla priorytetyzacji, przydziału zadań i estymacji kosztów, dla których rozdział przedstawia wyłącznie zaplanowaną metodykę. Opis procedury badawczej nie jest zatem równoznaczny z potwierdzeniem jej wykonania i miejsca te oznaczono osobno.

\section{Cele i zakres weryfikacji}

Weryfikacja rozwiązania obejmuje trzy uzupełniające się obszary. Pierwszym jest poprawność funkcjonalna, rozumiana jako zgodność operacji z wymaganiami i regułami modelu. Drugim jest poprawność techniczna, obejmująca granice modułów, kontrolę dostępu, obsługę błędów i trwałość danych. Trzecim jest jakość wspomagania decyzji.

W odniesieniu do ostatniego obszaru przyjęto następujące pytania badawcze:

\begin{enumerate}
    \item Czy ocena ryzyka wskazuje zadania zagrożone opóźnieniem skuteczniej niż prosta reguła oparta wyłącznie na terminie?
    \item Czy gromadzone dane uczące powstają poprawnie, to znaczy czy rejestrowane cechy odpowiadają stanowi z chwili prognozy, a etykiety rzeczywistym wynikom realizacji?
    \item Czy ranking backlogu odpowiada zadanym preferencjom i jak silnie zmienia się po modyfikacji wag lub danych?
    \item Czy propozycje przydziałów poprawiają rozdział pracy względem prostych metod bazowych?
    \item Czy estymacje modelu językowego są kompletne, spójne rachunkowo i dostatecznie stabilne?
    \item Czy czas oczekiwania i sposób prezentacji wyników pozwalają wykorzystać analizy podczas pracy z projektem?
\end{enumerate}

Odpowiedź na każde pytanie wymaga odpowiedniej podstawy porównania. Nie należy utożsamiać zgodności wyniku z zakodowanym wzorem z trafnością założeń tego wzoru.

Przykładowo test potwierdzający, że zadanie bez wykonawcy otrzymuje dodatkowe punkty ryzyka, sprawdza implementację reguły. Dopiero porównanie z rzeczywistymi rezultatami realizacji pozwala ocenić, czy reguła pomaga identyfikować przyszłe opóźnienia.

\section{Warunki odtwarzalności badań}

Przed wykonaniem eksperymentów należy utrwalić wersję aplikacji, konfigurację oraz dane wejściowe. W przypadku mechanizmów korzystających z daty bieżącej konieczne jest również ustalenie daty odniesienia.

Opis środowiska powinien zawierać:

\begin{itemize}
    \item identyfikator wersji kodu;
    \item wersje środowiska .NET, Node.js i PostgreSQL;
    \item parametry komputera lub środowiska uruchomieniowego;
    \item konfigurację procesorów zadań asynchronicznych;
    \item identyfikator i wersję zbioru danych;
    \item nazwę modelu językowego i konfigurację wywołań;
    \item nazwę modelu predykcji ryzyka zapisaną przy uruchomieniu analizy;
    \item datę odniesienia wykorzystywaną w obliczeniach.
\end{itemize}

Dla heurystyk deterministycznych te same dane, parametry i data odniesienia powinny prowadzić do tego samego wyniku. W przypadku modelu językowego powtarzalność należy sprawdzić przez niezależne wywołania.

Istotne jest uwzględnienie pamięci podręcznej estymacji kosztów. Wielokrotne pobranie tego samego zapisanego rezultatu nie jest badaniem stabilności modelu. Eksperyment powtarzalności musi obejmować rzeczywiste wywołania modelu z identycznym kontekstem, z pominięciem ponownego wykorzystania wcześniejszej estymacji.

% DO UZUPEŁNIENIA PO PRZEPROWADZENIU BADAŃ:
% wersja kodu, środowisko, konfiguracja, identyfikatory danych,
% liczba powtórzeń i terminy wykonania eksperymentów.

\section{Wyniki wykonania zestawu testów}
\label{sec:wyniki-testow}

Zautomatyzowany zestaw testów uruchomiono dla bieżącej wersji rozwiązania. Wykonano 45 testów, z których wszystkie zakończyły się powodzeniem; żaden nie został pominięty. Podział na projekty testowe przedstawia tabela~\ref{tab:wyniki-testow}.

\begin{table}[htbp]
    \centering
    \small
    \caption{Wynik uruchomienia zautomatyzowanego zestawu testów}
    \label{tab:wyniki-testow}
    \begin{tabularx}{\textwidth}{@{}Xrrr@{}}
        \toprule
        \textbf{Projekt testowy} & \textbf{Wykonane} & \textbf{Powodzenie} & \textbf{Niepowodzenie} \\
        \midrule
        Predykcja ryzyka (warstwa aplikacji) & 21 & 21 & 0 \\
        Testy architektoniczne & 9 & 9 & 0 \\
        Tożsamość i dostęp (warstwa aplikacji) & 8 & 8 & 0 \\
        Harmonogramowanie (warstwa infrastruktury) & 7 & 7 & 0 \\
        \midrule
        \textbf{Razem} & \textbf{45} & \textbf{45} & \textbf{0} \\
        \bottomrule
    \end{tabularx}

    \par\smallskip
    \begin{minipage}{\textwidth}
        \footnotesize
        Źródło: opracowanie własne na podstawie uruchomienia zestawu testów
        dla opisanej wersji rozwiązania.
    \end{minipage}
\end{table}

Rozkład ten sam w sobie stanowi wynik wart odnotowania. Najliczniejszy zestaw dotyczy modelu predykcyjnego, co odpowiada temu, że jest to element, którego niepoprawne działanie najtrudniej zauważyć: błąd w odwzorowaniu cech lub w kolejności przekształceń nie powoduje awarii, lecz ciche zwracanie wartości innych niż ocenione. Zakres pokrycia pozostaje natomiast nierównomierny, ponieważ moduły realizacji pracy, powiadomień i estymacji kosztów nie posiadają własnych projektów testowych.

Liczba wykonanych testów nie jest miarą jakości rozwiązania i nie należy jej tak interpretować. Wynik ten potwierdza wyłącznie, że sprawdzane własności zachodzą w badanej wersji kodu.

\section{Testy architektoniczne}

W repozytorium znajduje się projekt \texttt{PlanWise.ArchitectureTests}. Zawiera testy kierunku zależności dla wszystkich ośmiu modułów oraz test ograniczający bezpośrednie odwołania między modułami.

Weryfikowane są reguły strukturalne związane z podziałem na warstwy i izolacją obszarów funkcjonalnych. Ich celem jest wykrywanie zmian, które naruszają przyjętą organizację rozwiązania, na przykład wprowadzają niedozwoloną zależność między projektami.

Testy architektoniczne mają znaczenie szczególnie podczas rozbudowy wspólnych kontraktów. Umożliwiają sprawdzenie, czy dodanie nowego mechanizmu analitycznego nie prowadzi do bezpośredniego wykorzystania wewnętrznych elementów innego modułu.

Zakres tych testów jest jednak ograniczony. Nie potwierdzają one poprawności zapytań do bazy, skuteczności autoryzacji ani atomowości operacji obejmujących kilka kontekstów danych. Nie wykrywają również każdej zależności semantycznej wynikającej z zawartości kontraktów.

Wszystkie dziewięć testów architektonicznych zakończyło się powodzeniem w uruchomieniu opisanym w sekcji~\ref{sec:wyniki-testow}. Oznacza to, że rozbudowa modułów o granice modeli, kolejne implementacje oraz część odpowiedzialną za dane uczące nie naruszyła przyjętego kierunku zależności ani izolacji obszarów funkcjonalnych.

\section{Testy jednostkowe i weryfikacja reguł domenowych}

W module \texttt{IdentityAccess} znajdują się testy dotyczące obsługi sesji, wybranych procedur użytkownika i walidacji. Obejmują między innymi:

\begin{itemize}
    \item zapis skrótu tokenu odświeżania;
    \item unieważnienie poprzedniego tokenu podczas odnowienia sesji;
    \item odrzucenie tokenu wygasłego;
    \item unieważnienie aktywnego tokenu przy wylogowaniu;
    \item normalizację danych podczas rejestracji;
    \item odrzucenie zduplikowanego adresu e-mail;
    \item odrzucenie niepoprawnego hasła;
    \item zatrzymanie niepoprawnego żądania przed wykonaniem procedury obsługi.
\end{itemize}

Testy te sprawdzają wybrane zachowania w izolacji. Nie zastępują weryfikacji całego procesu uwierzytelniania z rzeczywistą bazą danych i klientem HTTP.

Dla funkcji projektowych istotne są reguły przejść stanów. Weryfikacja powinna objąć rozpoczęcie sprintu planowanego, zakończenie sprintu aktywnego i odrzucenie niedozwolonych przejść. Osobno należy sprawdzić zachowanie zadania przy dołączeniu do sprintu i usunięciu tego przypisania.

Testy częściowej aktualizacji powinny rozróżniać pominięcie pola, przekazanie nowej wartości i jawne wyczyszczenie wartości. Dotyczy to zwłaszcza wykonawcy, oszacowania, terminu i sprintu.

Dla algorytmów należy wybierać przypadki ujawniające istotne właściwości, a nie jedynie powtarzające implementację. Przykładem jest sprawdzenie, czy zwiększenie wartości biznesowej nie obniża oceny tego samego zadania przy niezmienionych pozostałych danych i wagach.

Oddzielnej kontroli wymagają wartości graniczne: brak zadań, brak oszacowań, zerowa pojemność, brak wykonawców i niepoprawny graf zależności.

\section{Testy integracyjne i scenariusze funkcjonalne}

Testy integracyjne, wymagające odrębnej bazy przygotowywanej w powtarzalny sposób, powinny sprawdzać współpracę punktów końcowych, procedur obsługi i bazy danych. Podstawowy scenariusz obejmuje utworzenie konta, projektu i zespołu, dodanie zadań, utworzenie sprintu i zmiany statusów, przy czym po każdej operacji należy zweryfikować zarówno odpowiedź API, jak i ponownie pobrany stan danych.

Szczególnie istotne są przypadki obejmujące więcej niż jeden moduł, w szczególności zastosowanie propozycji przydziałów wraz z niepowodzeniem w trakcie wykonywania kolejnych zmian. Pozwoli to określić, czy częściowa zmiana jest wykrywana i czy stan propozycji odpowiada rzeczywistym danym.

Dla zadań asynchronicznych weryfikacji wymaga zapis zlecenia, wybór właściwej procedury wykonawczej oraz zapis wyniku lub niepowodzenia, a osobno zachowanie po przerwaniu procesu i przy równoczesnym pobieraniu zleceń przez kilka instancji. Te ostatnie przypadki wyznaczają ograniczenia wdrożeniowe, ponieważ wspólny zapis zleceń w bazie nie zapewnia sam przez się bezpiecznego przejmowania pracy. Dla mechanizmu Outbox oczekiwania powinny odpowiadać obecnej implementacji, która nie ponawia wiadomości oznaczonych jako przetworzone z błędem.

\section{Testy interfejsu i kontroli dostępu}

Weryfikacja interfejsu powinna objąć tworzenie i edycję zadania, filtrowanie tablicy, zmianę sprintu, przeciąganie kart oraz przegląd wyników analiz. Dwa przypadki wymagają szczególnej uwagi. Po odrzuceniu zmiany przez serwer aktualizacja optymistyczna musi przywrócić widok do stanu zgodnego z danymi i umożliwić rozpoznanie niepowodzenia. Zmianę kolejności trzeba natomiast sprawdzić zarówno bez filtrów, jak i przy aktywnym filtrze sprintu, ponieważ indeks widoczny w przefiltrowanej kolumnie może mieć inne znaczenie niż indeks w pełnym zbiorze przetwarzanym przez backend. Propagację zmian przez SignalR należy zweryfikować w dwóch równocześnie otwartych sesjach, wraz z ponownym połączeniem po utracie komunikacji.

Testy kontroli dostępu wymagają co najmniej dwóch kont i dwóch projektów o różnych członkostwach. Dla zadań, sprintów, wyników analiz, propozycji, zleceń asynchronicznych i połączeń SignalR należy sprawdzić, czy użytkownik spoza projektu uzyska dostęp po podaniu znanego identyfikatora zasobu; ochrona listy projektów nie wystarcza, jeżeli bezpośredni adres pozwala ominąć ograniczenie. Osobno trzeba odróżnić dostęp do projektu od uprawnień do konkretnego działania oraz objąć scenariuszami token wygasły, token unieważniony i zmianę członkostwa podczas aktywnej sesji, obejmującą także wcześniej ustanowione połączenie czasu rzeczywistego.

\section{Projekt danych eksperymentalnych}

Do oceny mechanizmów analitycznych potrzebne są różne rodzaje danych. Małe, ręcznie przygotowane przypadki umożliwiają sprawdzenie obliczeń. Większe zbiory syntetyczne pozwalają badać reakcję na skalę i zmianę parametrów. Dane historyczne są natomiast potrzebne do oceny trafności prognoz względem rzeczywistych rezultatów.

Zbiór scenariuszy powinien uwzględniać:

\begin{itemize}
    \item projekty o małej i dużej liczbie zadań;
    \item zadania niezależne, łańcuchy i rozgałęzione grafy zależności;
    \item zróżnicowane wartości biznesowe i oszacowania;
    \item równomiernie i nierównomiernie obciążone zespoły;
    \item brakujące dane;
    \item niejednolite kompetencje i dostępności;
    \item zmiany zakresu, terminów i parametrów zespołu.
\end{itemize}

Dla danych syntetycznych należy zapisać reguły generowania i ziarno generatora. Scenariusze powinny zawierać również przypadki niekorzystne dla ocenianej metody.

Nie należy generować etykiet opóźnienia bezpośrednio z funkcji \texttt{RiskScorer}, a następnie używać ich jako dowodu trafności tej funkcji. Taki eksperyment potwierdzałby jedynie zgodność z przyjętą regułą generowania danych.

Podobnie ranking referencyjny nie powinien być wyznaczany przez te same wagi, które są oceniane, jeżeli celem jest badanie użyteczności priorytetyzacji.

Brak rzeczywistych danych historycznych ogranicza zakres możliwych wniosków. W takim przypadku można ocenić poprawność, wrażliwość i zachowanie w scenariuszach, ale nie rzeczywistą skuteczność prognozowania opóźnień.

Odrębnym źródłem danych jest zbiór obserwacji gromadzony przez system podczas kolejnych uruchomień analizy ryzyka. Różni się on od pozostałych tym, że powstaje wyłącznie wraz z upływem czasu i nie może zostać wygenerowany na potrzeby eksperymentu. Na etapie sporządzania niniejszego opisu jego liczebność nie pozwalała na dopasowanie modelu, dlatego przewidziano dla niego dwa odrębne rodzaje oceny, omówione w sekcji~\ref{subsec:ocena-danych-uczacych}: weryfikację poprawności samego mechanizmu rejestrowania, możliwą do wykonania natychmiast, oraz ocenę jakości modelu, możliwą dopiero po zgromadzeniu wystarczającej liczby obserwacji.

\section{Ocena mechanizmu ryzyka i prognoz sprintu}

\subsection{Ocena heurystyki ryzyka}

Pierwszy etap obejmuje sprawdzenie reakcji na czynniki wejściowe: terminy graniczne, liczbę blokad, próg rozmiaru, brak wykonawcy i ograniczenie końcowej wartości. Ocena zdolności przewidywania wymaga natomiast zadań obserwowanych przed upływem terminu, przy czym zadania już przeterminowane powinny tworzyć osobną grupę, ponieważ rozpoznanie istniejącego opóźnienia nie jest prognozą. Jako rozwiązanie bazowe można przyjąć regułę wskazującą zadania z terminem w ciągu trzech dni, oceniając obie metody dla tych samych zadań, momentów obserwacji i definicji opóźnienia.

Dla wybranego progu oceny można obliczyć precyzję i czułość:

\begin{align}
    \operatorname{Precision}
    &= \frac{TP}{TP+FP},\\
    \operatorname{Recall}
    &= \frac{TP}{TP+FN},
    \label{eq:testy-precision-recall}
\end{align}

gdzie \(TP\) oznacza poprawnie wskazane opóźnienia, \(FP\) fałszywe alarmy, a \(FN\) opóźnienia niewskazane przez metodę. Jeżeli mianownik jest równy zero, należy zgłosić brak możliwości wyznaczenia miary.

Przy ograniczonej liczbie zadań, które kierownik może przeanalizować, użyteczna jest również trafność wśród pierwszych \(k\) wskazań:

\begin{equation}
    \operatorname{Precision@k}
    =
    \frac{\text{liczba opóźnionych zadań w pierwszych }k\text{ wskazaniach}}{k}.
    \label{eq:testy-precision-k}
\end{equation}

Próg i wartość \(k\) należy ustalić przed oceną końcowego zbioru. Dobór ich na tych samych danych, na których raportowany jest wynik, prowadziłby do zawyżenia oceny.

Ponieważ obecna punktacja nie jest skalibrowanym prawdopodobieństwem, podstawowa ocena powinna dotyczyć jakości wskazań i rankingu. Nie należy interpretować samego zakresu od zera do jedności jako potwierdzenia probabilistycznego charakteru wyniku.

\subsection{Ocena prognoz sprintu}

Ocena prognoz sprintu stała się możliwa po wprowadzeniu jawnego przeliczenia dostępności zespołu na tempo realizacji. Należy przy tym raportować osobno przypadki, w których tempo zmierzono na zakończonych sprintach projektu, i te, w których przyjęto wartość założoną, ponieważ są to prognozy o różnej podstawie.

Mierzyć można błąd przewidywanej liczby wykonanych punktów oraz daty zakończenia. Dla \(n\) obserwacji średni błąd bezwzględny wynosi:

\begin{equation}
    \operatorname{MAE}
    =
    \frac{1}{n}\sum_{i=1}^{n}
    \left|\hat{y}_i-y_i\right|,
    \label{eq:testy-mae}
\end{equation}

gdzie \(\hat{y}_i\) jest prognozą, a \(y_i\) wynikiem rzeczywistym. Dla dat różnicę należy wyrazić w dniach.

Warianty zapisane pod nazwami P50 i P90 należy oceniać jako wariant bazowy i pesymistyczny. Ich interpretacja jako kwantyli wymagałaby odrębnego modelu i sprawdzenia częstości dotrzymywania wskazanych terminów.

\subsection{Ocena mechanizmu gromadzenia danych uczących}
\label{subsec:ocena-danych-uczacych}

Weryfikacja tego mechanizmu rozpada się na dwa etapy, dotyczące różnych przedmiotów oceny i wykonalne w różnym czasie.

Pierwszy etap obejmuje poprawność samego rejestrowania i jest możliwy niezależnie od liczby zgromadzonych obserwacji. Sprawdzeniu podlega zgodność zapisanych cech ze stanem zadania w dniu analizy, w szczególności znak liczby dni pozostałych do terminu dla zadania przeterminowanego, zapisanie dokładnie jednej obserwacji na oceniane zadanie w przebiegu oraz poprawność wszystkich czterech ścieżek etykietowania: nadania stanu wykluczającego etykietę zadaniu bez terminu, uzupełnienia wyniku po zakończeniu, oznaczenia jako cenzurowanej obserwacji dotyczącej zadania otwartego po terminie wraz z jej późniejszym rozstrzygnięciem oraz pozostawienia nierozstrzygniętą obserwacji, której zadanie usunięto lub pozbawiono terminu. Osobno należy potwierdzić, że rejestrowane są wielkości surowe: wykrycie wartości progowanych lub ważonych oznaczałoby przeniesienie założeń heurystyki do danych uczących jej następcy.

Drugi etap, dotyczący jakości modelu, wymaga zgromadzenia wystarczającej liczby rozstrzygniętych obserwacji. Poprzedzać go musi charakterystyka zbioru obejmująca liczbę obserwacji w podziale na stany wyniku, liczbę odrębnych zadań i projektów, udział obserwacji cenzurowanych oraz zakres czasowy gromadzenia.

Podział na część uczącą i testową musi grupować obserwacje według zadania oraz uwzględniać porządek czasowy, a zastosowaną procedurę należy podać wraz z wynikiem: podział losowy prowadziłby tu do oceny zawyżonej i zarazem trudnej do odróżnienia od poprawnej. Obserwacje cenzurowane wymagają przy tym odmiennego traktowania zależnie od sformułowania zadania, ponieważ w ujęciu klasyfikacyjnym informacja o wystąpieniu opóźnienia jest pełna, natomiast w regresyjnym zapisana liczba dni stanowi jedynie dolne ograniczenie.

Porównanie z modelem bazowym można przeprowadzić na tych samych obserwacjach testowych bez ponownego uruchamiania systemu, ponieważ wraz z każdą obserwacją zapisano prognozę z danego dnia oraz nazwę modelu, który ją sporządził. Zgromadzone dane pozwolą zatem ocenić również jakość obecnej heurystyki, a nie wyłącznie jej następcy.

\subsection{Wyniki oceny modelu uczonego}

Ocenę modelu predykcji opóźnień przeprowadzono na zewnętrznym zbiorze danych, ponieważ dane gromadzone przez system nie osiągnęły wymaganej liczebności. Zastosowaną metodykę, przyjętą definicję etykiety, wykryte źródła przecieku informacji oraz uzyskane wartości miar przedstawiono w sekcji~\ref{subsec:protokoly-oceny}, aby zachować je w bezpośrednim sąsiedztwie opisu samego modelu. W tym miejscu odniesiono je do przyjętych kryteriów oceny rozwiązania.

Wynik należy uznać za częściowe spełnienie kryterium jakości predykcji. Wykazano przewagę nad metodą odniesienia w protokole odpowiadającym rzeczywistym warunkom użycia, czyli na projektach nieobecnych w zbiorze uczącym, przy wartości pola pod krzywą ROC wynoszącej 0,589 wobec 0,500. Jednocześnie wielkość tej przewagi nie uprawnia do twierdzenia o wysokiej trafności prognoz pojedynczych zadań.

Nie wykazano natomiast przewagi modelu uczonego nad heurystyką stosowaną w systemie. Obie metody nie zostały dotąd porównane na tych samych danych, ponieważ heurystyka operuje na cechach częściowo niedostępnych w zbiorze zewnętrznym, przede wszystkim na terminie pojedynczego zadania. Porównanie takie jest wykonalne dopiero na danych gromadzonych przez wdrożenie i pozostaje zadaniem otwartym. Z tego powodu model uczony nie został uczyniony metodą domyślną.

Kryterium przejrzystości rekomendacji uznano za spełnione w stopniu wyższym niż dla metody bazowej, ponieważ wkłady poszczególnych cech sumują się dokładnie do wartości funkcji decyzyjnej, co podlega sprawdzeniu testem. Kryterium prezentowania ograniczeń uznano za spełnione przez dołączanie do każdego uruchomienia informacji o pochodzeniu danych uczących, o braku kalibracji na danych wdrożenia oraz o tym, która część wyniku nie pochodzi z modelu.

\subsection{Pomocnicze studium czasu obsługi zgłoszenia}

Odrębnie przeprowadzono studium na zbiorze odrzuconym na etapie doboru danych, opisanym w sekcji~\ref{subsec:dobor-zbioru-danych}. Ponieważ zbiór ten nie pozwalał zbudować etykiety opóźnienia, zmieniono pytanie badawcze na przewidywanie czasu obsługi zgłoszenia liczonego od jego utworzenia do rozwiązania.

Studium to zachowano w pracy przede wszystkim ze względu na wykrytą w nim pułapkę metodyczną, która ma znaczenie ogólniejsze. Eksport danych stanowi migawkę wykonaną w ustalonym dniu, wobec czego zgłoszenie utworzone niedawno może wyglądać na rozwiązane szybko wyłącznie dlatego, że inaczej nie zdążyłoby się jeszcze rozwiązać. Pierwsza wersja eksperymentu, ucząca się na zgłoszeniach rozwiązanych i dzieląca dane według daty, dała wartość pola pod krzywą ROC równą 0,385, czyli gorszą od losowej. Mierzyła bowiem artefakt obcięcia obserwacji, a nie zdolność predykcyjną.

Zastosowano dwie poprawki. Wprowadzono okno obserwacji, dopuszczające wyłącznie zgłoszenia utworzone co najmniej rok przed datą migawki, tak aby każde miało tyle samo czasu na rozwiązanie. Zgłoszenia nierozwiązane potraktowano ponadto jako obserwacje klasy negatywnej, zamiast je odrzucać; ich pominięcie stanowiłoby selekcję obserwacji przetrwałych, ponieważ zgłoszenia nigdy nierozwiązane są systematycznie tymi najtrudniejszymi. Po tych poprawkach udziały klas w zbiorze uczącym i testowym zrównały się z dokładnością do około jednego punktu procentowego.

Wyniki dla progu trzydziestu dni przedstawia tabela~\ref{tab:wyniki-studium-pomocniczego}.

\begin{table}[htbp]
    \centering
    \small
    \caption{Przewidywanie rozwiązania zgłoszenia w ciągu trzydziestu dni}
    \label{tab:wyniki-studium-pomocniczego}
    \begin{tabularx}{\textwidth}{@{}Xll@{}}
        \toprule
        \textbf{Zestaw cech} & \textbf{ROC AUC} & \textbf{Przyrost w pierwszym decylu} \\
        \midrule
        Metoda odniesienia & 0,500 & 0,89 \\
        Metadane & 0,741 & 1,25 \\
        Tekstowe & 0,647 & 3,75 \\
        Łączone & \textbf{0,795} & \textbf{4,06} \\
        \bottomrule
    \end{tabularx}

    \par\smallskip
    \begin{minipage}{\textwidth}
        \footnotesize
        Źródło: opracowanie własne. Zbiór testowy liczył 7\,968 zgłoszeń
        przy udziale klasy pozytywnej 4,9\%.
    \end{minipage}
\end{table}

Wyniki te są wyraźnie lepsze niż w zadaniu predykcji opóźnienia sprintu, czego nie należy jednak odczytywać jako sprzeczności. Są to dwa różne zadania, o różnych definicjach zdarzenia i różnych udziałach klas, a łatwiejsze z nich dotyczy kolejki zgłoszeń publicznych, nie zaś dotrzymywania zobowiązań sprintu.

Dla PlanWise istotne są dwa wnioski. Pierwszy dotyczy komplementarności: żaden zestaw cech osobno nie osiąga wyniku zestawu łączonego, a znaczenie cech tekstowych rośnie wraz z wydłużaniem horyzontu prognozy, podczas gdy same metadane tracą wówczas zdolność różnicującą. Drugi powtarza obserwację z zadania głównego: także tutaj model o dobrym uporządkowaniu wypada gorzej od modelu stałego pod względem miary Briera, co potwierdza, że kalibracja wymaga osobnego kroku i nie wynika z jakości rankingu.

Ze studium tego nie wyprowadzono natomiast wniosków o wielkości nakładu pracy dla zadań w PlanWise. Mediana czasu obsługi w badanym zbiorze wynosi 128 dni, co odpowiada publicznej kolejce zgłoszeń, a nie realizacji zadań w sprincie.

\section{Ocena priorytetyzacji backlogu}

Ocena priorytetyzacji powinna obejmować poprawność funkcji, stabilność rankingu i zgodność z niezależnym punktem odniesienia. Jako metody bazowe można wykorzystać bieżącą kolejność backlogu oraz sortowanie wyłącznie według wartości biznesowej, przy zachowaniu tego samego zakresu zadań.

Analiza wrażliwości polega na zmianie wag i ponownym wyznaczeniu kolejności. Scenariusze uzupełniające powinny obejmować dodanie zadania o skrajnej wartości, uzupełnienie brakującego oszacowania i aktualizację ocen ryzyka, co pozwoli określić, czy niewielka zmiana danych prowadzi do dużej reorganizacji backlogu.

Przeciętną zmianę pozycji można opisać jako:

\begin{equation}
    D_{\mathrm{rank}}
    =
    \frac{1}{n}\sum_{i=1}^{n}
    \left|r_i^{(a)}-r_i^{(b)}\right|,
    \label{eq:testy-zmiana-rankingu}
\end{equation}

gdzie \(r_i^{(a)}\) i \(r_i^{(b)}\) oznaczają pozycje tego samego zadania w dwóch rankingach.

Oceniający z doświadczeniem projektowym mogą przygotować niezależną kolejność na podstawie wspólnego opisu celu i ograniczeń, przy czym raportować należy również rozbieżności między nimi, ponieważ nie zawsze istnieje jeden poprawny ranking. Zgodność z oceną eksperta nie dowodzi poprawy wyniku projektu, lecz wskazuje, czy rekomendacja jest zrozumiała i odpowiada zadanym preferencjom.

Obecność dwóch implementacji wymaga rozszerzenia planu oceny o porównanie między nimi, obejmujące zgodność obu kolejności, ich zgodność z punktem odniesienia oraz kompletność wyniku, ponieważ model językowy może zwrócić mniej pozycji, niż otrzymał. Osobno należy zmierzyć powtarzalność, która ma dla obu metod odmienny charakter: heurystyka jest deterministyczna, natomiast dla modelu językowego trzeba wykonać wielokrotne wywołania z tym samym wejściem i opisać rozrzut pozycji, przyjmując za miarę wzór~\eqref{eq:testy-zmiana-rankingu} zastosowany do dwóch przebiegów tej samej metody.

Wynik tego porównania nie jest przesądzony. Model językowy operuje na opisach zadań, których ocena punktowa w ogóle nie analizuje, lecz czyni to kosztem determinizmu i możliwości prześledzenia obliczeń.

\section{Ocena harmonogramu i przydziałów}

\subsection{Weryfikacja obliczeń harmonogramu}

Dla małych grafów należy przygotować ręcznie wyznaczone terminy i zapasy czasu, obejmując przypadkami pojedyncze zadanie, łańcuch, rozgałęzienie i kilka zadań końcowych, a osobno ręczne nadpisania, brak oszacowania, graf cykliczny i zmianę daty odniesienia. Wynik dla grafu zawierającego cykl powinien zostać oznaczony jako niewiarygodny, nawet jeżeli aplikacja potrafi go wyświetlić. Zgodność z CPM nie potwierdza przy tym wykonalności zasobowej: graf może dopuszczać równoległą realizację dwóch zadań wymagających tego samego wykonawcy.

\subsection{Porównanie metod przydzielania}

Obie zaimplementowane metody należy porównać ze sobą oraz z metodami prostszymi: cyklicznym przydzielaniem kolejnych zadań i wyborem najmniejszego obciążenia względem pojemności bez preferencji kompetencji. Każda z nich musi otrzymać ten sam zbiór nieprzypisanych zadań i ten sam początkowy stan zespołu, a istniejące przydziały powinny pozostać niezmienione we wszystkich wariantach.

Dobór miary uległ przy tym zasadniczej zmianie. Dopóki jedyną metodą była heurystyka zachłanna, ocena musiała ograniczać się do równomierności obciążenia, mierzonej rozstępem względnego obciążenia \(u_j=L_j/C_j\), a wnioskowanie z niej o skróceniu projektu byłoby nieuprawnione. Solver optymalizuje natomiast bezpośrednio długość projektu, wobec czego to ona stanowi właściwą miarę porównawczą, a rozstęp obciążenia staje się wielkością pomocniczą. Uzupełniająco należy podać liczbę przydzielonych zadań, udział przydziałów z dopasowaniem kompetencji i liczbę naruszeń zdefiniowanych ograniczeń, przy czym samo dopasowanie warto zweryfikować względem ręcznie opisanych wymagań zadań, ponieważ zgodność fragmentu tytułu z nazwą kompetencji jest właściwością heurystyki, a nie miarą merytorycznej trafności wyboru.

Porównanie obu metod musi przy tym odbywać się względem tej samej wielkości i przy tych samych ograniczeniach. Samo zestawienie długości projektu wyznaczonej przez solver z sumą punktów przypisanych przez heurystykę nie byłoby porównaniem, ponieważ heurystyka nie wyznacza terminów. Poprawne postępowanie polega na rozplanowaniu wyborów heurystyki w tym samym modelu czasu, którym posługuje się solver, i porównaniu wynikających stąd terminów zakończenia. Taki sposób obliczania przewidywanej korzyści zastosowano w implementacji.

Ocena powinna ponadto rozróżniać przypadki zakończone rozwiązaniem dowiedzionym jako optymalne od tych, w których znaleziono jedynie rozwiązanie dopuszczalne, oraz osobno raportować udział uruchomień zakończonych wykonaniem metody zastępczej. Rosnący udział tych ostatnich wraz z wielkością projektu wyznacza praktyczną granicę stosowalności solvera i powinien zostać zmierzony dla projektów o różnej liczbie zadań.

Osobnej weryfikacji wymagają ograniczenia zadeklarowane jako dotrzymane. Należy sprawdzić na wygenerowanych propozycjach, czy istotnie żadne zadanie nie rozpoczyna się przed zakończeniem swoich poprzedników i czy żadnej osobie nie przypisano dwóch zadań o nakładających się przedziałach. Deklaracja dotrzymania ograniczenia jest bowiem elementem wyjaśnienia, a nie dowodem jego spełnienia.

\section{Ocena estymacji kosztów}

\subsection{Kompletność i spójność}

Pierwszy etap oceny dotyczy odpowiedzi jako danych aplikacyjnych. Należy sprawdzić możliwość jej odczytania, kompletność scenariuszy, zgodność ról i stawek oraz nieujemność godzin i kosztów.

Dla każdej pozycji pracy należy ponownie obliczyć iloczyn godzin i stawki. Osobno trzeba zweryfikować sumy oraz powiązanie pozycji kosztowych z poszczególnymi scenariuszami.

Udział odpowiedzi spełniających wszystkie zdefiniowane warunki można wyrazić jako:

\begin{equation}
    V =
    \frac{n_{\mathrm{valid}}}{n_{\mathrm{all}}},
    \label{eq:testy-poprawnosc-estymacji}
\end{equation}

gdzie \(n_{\mathrm{valid}}\) oznacza liczbę poprawnych odpowiedzi, a \(n_{\mathrm{all}}\) wszystkie próby. Należy raportować wynik pierwszej próby oraz wynik po zastosowaniu dopuszczonego ponowienia.

\subsection{Powtarzalność i trafność}

Dla niezmienionego kontekstu należy wykonać ustaloną wcześniej liczbę niezależnych wywołań. Porównanie powinno obejmować koszt scenariusza realistycznego, rozkład godzin między role, założenia i rekomendacje.

Względne zróżnicowanie kosztu można opisać współczynnikiem zmienności:

\begin{equation}
    CV=\frac{s_K}{\overline{K}},
    \label{eq:testy-cv}
\end{equation}

gdzie \(s_K\) oznacza odchylenie standardowe, a \(\overline{K}>0\) średnią kosztów. Niska zmienność oznacza stabilność, lecz nie potwierdza trafności.

Jeżeli dostępne są rzeczywiste koszty porównywalnego zakresu, można obliczyć względny błąd:

\begin{equation}
    E_K =
    \frac{|\hat{K}-K|}{K},
    \qquad K>0.
    \label{eq:testy-blad-kosztu}
\end{equation}

Zakres i kategorie kosztów muszą być zgodne. Estymacji pozostałych prac nie należy porównywać bezpośrednio z całkowitym wydatkiem projektu obejmującym również prace zakończone.

Przy braku danych rzeczywistych możliwe jest porównanie z niezależną estymacją ekspercką lub szczegółową kalkulacją scenariusza. Wynik należy wtedy określać jako zgodność z przyjętym odniesieniem, a nie trafność względem kosztu rzeczywistego.

\subsection{Ocena rekomendacji redukcji}

Rekomendacje należy oceniać pod względem związku z backlogiem, wykonalności, wpływu na cel projektu i wiarygodności oszczędności, sprawdzając przy tym, czy propozycje nie dublują tego samego efektu. Zapisanie wyboru redukcji w aplikacji nie jest dowodem uzyskania oszczędności; potwierdzenie efektu wymaga realizacji zmiany i porównania kosztów w zgodnym zakresie.

\section{Ocena wydajności i odporności na błędy}

Pomiary wydajności powinny rozróżniać zwykłe żądania API od analiz asynchronicznych, a dla tych ostatnich osobno ujmować czas przyjęcia zlecenia, oczekiwania w kolejce, obliczeń i udostępnienia wyniku. Dla mechanizmów korzystających z modelu językowego konieczne jest dodatkowo oddzielenie czasu komunikacji z usługą zewnętrzną od czasu przetwarzania lokalnego, a dla solvera — zmierzenie udziału uruchomień kończących się metodą zastępczą. Wyniki należy raportować dla określonego rozmiaru projektu, podając obok średniej medianę, percentyl 95 i odsetek błędów, przy progach akceptacji ustalonych przed analizą.

Scenariusze odporności powinny obejmować niedostępność modelu, niepoprawną odpowiedź, błąd zapisu i przerwanie wykonania zadania. Weryfikacji podlega, czy użytkownik otrzymuje jednoznaczną informację, czy wcześniejsze poprawne dane pozostają dostępne oraz czy wynik aktualny daje się odróżnić od starszego — prezentowanie poprzedniego rezultatu bez oznaczenia jego daty prowadziłoby do błędnej interpretacji.

\section{Prezentacja wyników i zagrożenia dla ich wiarygodności}

Wyniki testów powinny zawierać liczbę przypadków wykonanych, zakończonych poprawnie, zakończonych błędem i pominiętych, wraz ze wskazaniem wersji aplikacji. Dla eksperymentów porównawczych należy podać dane wejściowe, metodę bazową, miary, liczbę powtórzeń i rozrzut wyników, uwzględniając także scenariusze, w których PlanWise wypadł gorzej lub nie przygotował poprawnej rekomendacji.

% WYKONANE:
% - Raport testów automatycznych (sekcja "Wyniki wykonania zestawu testów").
% - Ocena modelu predykcji opóźnień na danych zewnętrznych
%   (wyniki w rozdziale 6, odniesienie do kryteriów w tym rozdziale).
% - Pomocnicze studium czasu obsługi zgłoszenia.
%
% DO UZUPEŁNIENIA PO WYKONANIU:
% 1. Wyniki scenariuszy funkcjonalnych i kontroli dostępu.
% 2. Weryfikacja mechanizmu rejestrowania danych uczących
%    oraz charakterystyka zgromadzonego zbioru.
% 3. Porównanie modelu uczonego z heurystyką na danych wdrożenia.
% 4. Analiza wrażliwości rankingu i porównanie obu metod priorytetyzacji.
% 5. Porównanie metod przydziału względem długości projektu.
% 6. Spójność, zmienność i błędy estymacji kosztów.
% 7. Pomiary czasu odpowiedzi i analizy niepowodzeń.

Podstawowym zagrożeniem dla wiarygodności jest wykorzystanie danych syntetycznych odzwierciedlających założenia ocenianego algorytmu. Kolejnym jest mała liczba projektów lub ocen ekspertów, ograniczająca możliwość uogólnienia.

W analizach historycznych należy uwzględnić zmiany terminów i zakresu oraz wykluczyć wykorzystanie informacji niedostępnych w chwili prognozy. W badaniu modelu językowego znaczenie mają wersja modelu, konfiguracja i sposób przygotowania kontekstu.

Wspomaganie decyzji można dodatkowo ocenić w badaniu użytkowników, mierząc czas wykonania zadania decyzyjnego, liczbę przeoczonych problemów i odsetek przyjętych rekomendacji. Sama deklaracja zadowolenia nie zastępuje jednak oceny poprawności decyzji.

Podsumowując stan weryfikacji: dla predykcji opóźnień dysponujemy wynikiem pomiaru wykazującym przewagę nad metodą odniesienia, przy jednoczesnym braku porównania z heurystyką stosowaną w systemie. Dla pozostałych mechanizmów rozdział przedstawia metodykę bez wyników. Analiza kodu i wykonany zestaw testów pozwalają wskazać zakres implementacji oraz jej ograniczenia, lecz nie stanowią potwierdzenia przewagi systemu nad metodami bazowymi w pozostałych obszarach.

Weryfikacja tezy pracy wymaga zatem uzupełnienia rozdziału o wyniki pozostałych badań. Osiągnięty dotąd rezultat pozwala natomiast zawęzić pytanie: nie chodzi już o to, czy model uczony jest w tej dziedzinie w ogóle możliwy do zbudowania, lecz o to, czy w warunkach konkretnego zespołu przewyższa dostrojoną heurystykę na tyle, by uzasadnić jego użycie.